\section{Introduction}
\label{sec:tpc31-introduction}

The fixed-shift prime-pair problem separates local congruence
compatibility from global signed cancellation.  Sieve methods can
encode the local constraints with great precision, but a two-target
lower bound encounters the parity barrier long before those local
constraints become the whole problem.  The present paper does not
claim to bypass that barrier \citep{HeathBrown1982Parity}.  It continues a narrower program: make
the unresolved M\"obius-tail correlation successively more rigid,
prove power saving on every sector accessible to exact incidence
geometry, and state the remaining signed input without disguising it
as a consequence of the absolute estimates.

TPC-18 reduced a failed slice to a primitive off-diagonal two-row
correlation with rows $m=\ell d$, orbit length $J$, row scale $Q$, and
nonzero affine determinant $h(m-n)$ \citep{WangTPC18}.  TPC-19 and
TPC-20 gave exact determinant-frequency and matched-channel normal
forms, while also identifying coefficient and zero-frequency gates
that the formal transforms do not estimate \citep{WangTPC19,WangTPC20}.
TPC-28 closed the selected low-conductor packet before centering;
TPC-29 removed large reflected-content and a content-rich sparse
wedge; and TPC-30 used the full gcd of the two targets to prune the
large-target-content sectors of the three retained raw channels
\citep{WangTPC28,WangTPC29,WangTPC30}.

The last of these steps leaves
\begin{equation}
 \cG_{m,n}(j):=(mj+h,nj+h)\le X^\eta,
 \label{eq:tpc31-intro-small-content}
\end{equation}
with the fully coprime target sector as the sharp model.  The proposed
next decomposition was by translated determinant windows and
determinant residue classes.  There is, however, a definitional issue
that must be settled before such a decomposition is legitimate.

If $r\mid mj+h$ and $s\mid nj+h$ are selected divisor variables, then
$g=(r,s)$ depends on that selection.  The quotient
\[
 \frac{m-n}{g}
\]
is useful inside the expanded divisor channel, but it is not a label
of the unexpanded raw kernel: the same row--orbit triple has many
factorizations and therefore potentially many values of $g$.  A
cell decomposition made with this quotient cannot be inserted before
the divisor sum without recording its multiplicity.

The factorization-independent replacement is
\begin{equation}
 c=\cG_{m,n}(j),\qquad
 \Delta^\#_{m,n}(j)=\frac{m-n}{c}.
 \label{eq:tpc31-intro-canonical-pair}
\end{equation}
On the primitive orbit, TPC-30's Euclidean identity gives
$c\mid m-n$, so $\Delta^\#$ is integral.  More is true.  Dividing the
two targets by their full gcd produces coprime integers
\[
 U=\frac{mj+h}{c},\qquad V=\frac{nj+h}{c},
\]
and direct subtraction gives
\begin{equation}
 \boxed{mV-nU=h\Delta^\#.}
 \label{eq:tpc31-intro-reduced-determinant}
\end{equation}
Thus $h\Delta^\#$ is the exact determinant of the canonically reduced
target pair.  This is the invariant used throughout the paper.

\subsection{The new cellwise estimate}

Fix an exact content $c$ and a set $\Omega_c$ of allowed normalized
determinants.  Since the physical row integers are distinct, for a
fixed left row $m$ and a fixed $\delta\in\Omega_c$ there is at most
one right row,
\[
 n=m-c\delta.
\]
The same statement holds with left and right reversed.  Hence the
corresponding bipartite row graph has both degrees at most the number
of allowed $\delta$ values in the physical difference range.  The
Schur test and the two physical row $\ell^2$ bounds turn this degree
bound into a weighted row estimate.  Exact content $c$ also confines
the orbit variable to one residue class modulo $c$, giving the factor
$1+J/c$.

This proves the master estimate
\begin{equation}
 \left|\cS\bigl(c\le Z,\Delta^\#\in\Omega_c\bigr)\right|
 \ll X^\eps Q\log(2L)
 \sum_{c\le Z}\left(1+\frac Jc\right)
 \cH_c(\Omega_c),
 \label{eq:tpc31-intro-master}
\end{equation}
where $\cH_c$ is the number of allowed normalized determinants in the
physical range.  The estimate is before divisor-variable expansion
and therefore applies to each of the two mixed raw channels and the
both-ultra raw channel.

Two consequences are especially transparent.  For arbitrary
translated intervals $I_c$ of length at most $W$,
\begin{equation}
 \left|\cS(c\le Z,\Delta^\#\in I_c)\right|
 \ll X^\eps Q(W+1)\bigl(Z+J\log(2Z)\bigr).
 \label{eq:tpc31-intro-window}
\end{equation}
For one residue class $a_c\pmod q$ at each content,
\begin{equation}
 \left|\cS(c\le Z,\Delta^\#\equiv a_c\!\pmod q)\right|
 \ll X^\eps XQ
 \left(\frac ZX+\frac1Q+\frac1q\right).
 \label{eq:tpc31-intro-residue}
\end{equation}
The first estimate is translation-invariant: its interval may be
centered at $|\Delta^\#|\asymp Q/c$, far from the deleted near-row
sector.

\subsection{Removing the small-content restriction}

The cell estimate and TPC-30 are complementary.  Given a threshold
$C$, the part with $c>C$ is bounded by the large-content theorem,
whereas \eqref{eq:tpc31-intro-master} treats $c\le C$.  For a single
translated interval $I$ of length $W$ and a residue class $a\pmod q$
this gives
\begin{align}
 &\left|\cS(\Delta^\#\in I,
 \Delta^\#\equiv a\!\pmod q)\right|
 \notag\\
 &\quad\ll X^\eps XQ\left[
 \frac1J+\frac1C
 +\left(1+\frac Wq\right)
 \left(\frac CX+\frac1Q\right)
 \right].
 \label{eq:tpc31-intro-splice}
\end{align}
This is an all-content theorem: no condition such as
\eqref{eq:tpc31-intro-small-content} remains in the event being
estimated.  Optimization of $C$ gives a fixed power whenever
$1+W/q$ is polynomially smaller than $Q$.

For a residue class over the full normalized-determinant range, the
dependence on $c$ improves the count and yields
\begin{equation}
 \left|\cS(\Delta^\#\equiv a\!\pmod q)\right|
 \ll X^\eps XQ
 \left(\frac1J+\frac1C+\frac CX+\frac1Q+\frac1q\right).
 \label{eq:tpc31-intro-full-residue}
\end{equation}
Every fixed exponent below
$\min\{1-\beta,\beta,\rho\}$ is therefore available when
$q=X^{\rho+o(1)}$.

\subsection{Why the theorem stops at one cell}

The gain in \eqref{eq:tpc31-intro-master} is cardinality gain.  If a
determinant range is partitioned into windows and every window is
estimated absolutely, their cardinalities sum back to the size of the
full range.  Likewise, summing the absolute majorants over all
$a\pmod q$ exactly restores all row pairs in the layer $c\mid m-n$.
There is no hidden global saving.

For the signed physical amplitudes, the discrete Fourier transform in
$a$ is exact.  Its $r=0$ coefficient is the unweighted sum over all
residue cells, namely the original exact-content contribution.  This
is an auxiliary residue-label transform, not the Poisson zero frequency
in the orbit variable and not TPC-20's centered-kernel zero mode.
Bounds for the nonzero frequencies or for equidistribution of the
cells do not, by themselves, bound that mean.  A candidate next route
must construct a factorable representation of the prime--M\"obius
coefficient while retaining the fixed residue factors, pair mask, and
matched channel structure; it cannot be obtained by merely repeating
the present Schur estimate with a Fourier notation.

\subsection{Scope}

All estimates here are unconditional consequences of the inherited
physical packet, TPC-30's proved large-content bound, and elementary
finite-dimensional analysis.  They do not use a Chowla--Elliott
hypothesis, a prime-pair conjecture, or an unverified transfer of a
Kloosterman or spectral large-sieve theorem.  Fixed-form logarithmic
two-point cancellation remains an important benchmark
\citep{Tao2016}, but it does not provide the required uniformity in
the polynomially growing row, content, cutoff, and mask families.

\Cref{sec:tpc31-interface} states the physical interface.
\Cref{sec:tpc31-canonical} proves the canonical reduction.
\Cref{sec:tpc31-row-entropy} gives the row-graph theorem, and
\cref{sec:tpc31-small-transfer} transfers it to the three raw
channels.  The all-content splice is in
\cref{sec:tpc31-all-content}; the high-$\beta$ ledger is in
\cref{sec:tpc31-high-beta}.  The exact Fourier and reassembly boundary
is isolated in \cref{sec:tpc31-zero-mode}.  The final sections compare
the result with the preceding regions and record the finite
certificate.
